% Generated draft: requires mathematical and source review.
\documentclass[11pt]{article}
\usepackage[margin=25mm]{geometry}
\usepackage{fontspec}
\setmainfont{Latin Modern Roman}
\usepackage{amsmath,amssymb,mathtools,unicode-math}
\setmathfont{Latin Modern Math}
\usepackage{enumitem,needspace,xurl,hyperref}
\setlength{\parindent}{0pt}
\setlength{\parskip}{5pt}
\setlength{\emergencystretch}{3em}
\newcommand{\N}{\mathbb{N}}
\newcommand{\Z}{\mathbb{Z}}
\newcommand{\Q}{\mathbb{Q}}
\newcommand{\R}{\mathbb{R}}
\newcommand{\C}{\mathbb{C}}
\newcommand{\F}{\mathbb{F}}
\newcommand{\A}{\mathbb{A}}
\newcommand{\PP}{\mathbb{P}}
\newcommand{\E}{\mathbb{E}}
\newcommand{\eps}{\varepsilon}
\newcommand{\dd}{\,\mathrm{d}}
\begin{document}
\section*{UnsolvedMath ID: 30000038}
\textbf{Source problem code:} catalogue entry
\subsection{Definitions and mathematical statement}
For a finite family $\Phi=(\phi_1,\ldots,\phi_N)$ of complex-valued functions in $L^p(\R^d)$, a conventional norm-closed shift-generated space, when $1\leq p<\infty$, is
\[
 S_p(\Phi)=\overline{\operatorname{span}\{\phi_j(\,\cdot-k):
                 1\leq j\leq N,\ k\in\Z^d\}}^{\,L^p}.
\]
Such a space is finitely generated and invariant under integer translations. For $A=(a_{ij})\in\C^{m\times N}$ define the linear combinations
\[
 \psi_i=\sum_{j=1}^N a_{ij}\phi_j\quad(1\leq i\leq m),
 \qquad \Psi=A\Phi.
\]
The source asks for an extension to $L^p(\R^d)$, $p\ne2$, of an $L^2$ theorem about generic linear combinations of generators. It does not state that theorem, the allowed number $m$, or what property the combinations must preserve.

For example, if the intended property is generation, define the length of a nonzero finitely generated space $V$ to be the minimum number $\ell_p(V)$ of generators. One possible precise target would ask, for appropriate $\ell_p(V)\leq m\leq N$, whether
\[
 S_p(A\Phi)=S_p(\Phi)=V
 \quad\text{for Lebesgue-almost every }A\in\C^{m\times N}.
\]
This is an explicitly conditional interpretation, not a recovered statement of the missing theorem. Preservation of frame bounds or stable translates is stronger than preservation of the closed span. Likewise ``generic'' could mean full measure, a residual set, or a specified algebraic condition; the source does not choose one. The displayed norm-closure convention does not decide the additional issues for $p=\infty$ or quasi-Banach $0<p<1$, and the original $p$ range is unspecified.

The supplied literature triage checked 21 August 2026 calls the general $L^p$ extension open while recording substantial results for $p=1$ and specialized shift-invariant settings. Thus the record must not be read as claiming the entire area is untouched, nor as supplying a full theorem for every $p\ne2$.
\subsection{Short English statement}
Extend a generic-linear-combination result for finitely generated $L^2$ shift-invariant spaces to other $L^p$ spaces. The missing $L^2$ theorem, meaning of genericity and preservation property prevent a unique precise target. The supplied assessment records partial $p=1$ and specialized results.
\subsection{Sources}
\begin{itemize}
\item Local UnsolvedMath registry, permanent ID 30000038, exact source code catalogue entry; $L^p$ Extensions for Shift-Invariant Spaces. Proposed year 2004, status partially\_solved; corrected\_verified extraction restores function-space notation and $p\ne2$.
\item \emph{Mini-Workshop: Wavelets and Frames (2004)}, Oberwolfach Reports, DOI 10.4171/owr/2004/10, pp. 479--544, \url{https://doi.org/10.4171/owr/2004/10}.
\item Supplied literature triage checked 2026-08-21: ``Oberwolfach report'', \url{https://doi.org/10.4171/owr/2004/10}; ``Shift-invariant spaces in $L^1$'', \url{https://doi.org/10.1006/jath.2002.3711}; ``Related $L^p$ work'', \url{https://doi.org/10.1016/j.jmaa.2005.05.055}. The dated OWR-derived triage requires independent verification of time-sensitive status and cited claims. Its description of later work is retained as supplied, without independently resolving publication chronology.
\end{itemize}
\end{document}
